\subsection{弗雷德霍姆映射的指标}

设 E、F 与 G 是局部凸空间。

设 \(u \colon \mathrm{E} \to \mathrm{F}\) 是弗雷德霍姆算子。向量空间 \(\operatorname{Ker}(u)\) 与 \(\operatorname{Coker}(u)\) 是有限维的（III, p. 41 的引理 1）。回忆整数

\[
\dim \operatorname{Coker} (u) - \dim \operatorname{Ker} (u) = \operatorname{codim} _ {\mathrm{F}} \operatorname{Im} (u) - \dim \operatorname{Ker} (u) \tag {1}
\]

称为 u 的指标，并记作 ind(u)（A, V, p. 126）。

如果 \(u \colon \mathrm{E} \to \mathrm{F}\) 与 \(v \colon \mathrm{F} \to \mathrm{G}\) 是弗雷德霍姆映射，那么 \(v \circ u\) 也是（III, p. 41, 注记 2），并且我们有（A, V, p. 127, 引理 2）

\[
\operatorname{ind} (v \circ u) = \operatorname{ind} (v) + \operatorname{ind} (u).\tag{2}
\]

假设 E 与 F 是分离的，并设 \(u\colon\mathrm{E}\to\mathrm{F}\) 是弗雷德霍姆算子；采用 III, p. 42 的 Prop. 2 的条件 (iv) 的记号。我们于是有 \(\mathrm{ind}(u)=\mathrm{dim}(\mathrm{F}_{2})-\mathrm{dim}(\mathrm{E}_{2})\)。给这些空间中每一个的对偶赋以弱拓扑（相应地，紧开拓扑、有界收敛拓扑）。那么 \(\mathrm{E}^{\prime}\) 等同于 \(\mathrm{E}_{1}^{\prime}\) 与 \(\mathrm{E}_{2}^{\prime}\) 的拓扑直和，\(\mathrm{F}^{\prime}\) 等同于 \(\mathrm{F}_{1}^{\prime}\) 与 \(\mathrm{F}_{2}^{\prime}\) 的拓扑直和，且 \(^{t}u\) 诱导了从 \(\mathrm{F}_{1}^{\prime}\) 到 \(\mathrm{E}_{1}^{\prime}\) 上的同构并在 \(\mathrm{F}_{2}^{\prime}\) 上为零。于是转置 \(^{t}u\colon\mathrm{F}^{\prime}\to\mathrm{E}^{\prime}\) 是弗雷德霍姆映射（loc. cit.）。 \(^{t}u\) 的核是 \(\mathrm{F}_{2}^{\prime}\)，它的维数等于 \(\mathrm{F}_{2}\) 的维数，也就是 \(u\) 的余核的维数。从而

\[
\operatorname{ind} (u) = \dim \operatorname{Ker} (^ {t} u) - \dim \operatorname{Ker} (u).\tag{3}
\]

此外，\(^{t}u\) 的像是 \(E_{1}^{\prime}\)，因而 \(^{t}u\) 的余核的维数等于 \(E_{2}^{\prime}\) 的维数，也就是 u 的核 \(E_{2}\) 的维数。由此我们推出

\[
\operatorname{ind} (^ {t} u) = - \operatorname{ind} (u).\tag{4}
\]

假设 E 与 F 是分离的，并设 \(u\colon\mathrm{E}\to\mathrm{F}\) 是指标为 0 的弗雷德霍姆算子。那么由 III, p. 42 的 prop. 2，\(u\) 是严格态射。由于 \(\dim\operatorname{Ker}(u)=\operatorname{codim}_{\mathrm{F}}\operatorname{Im}(u)\)，映射 \(u\) 只要是单射或满射就是同构。

命题 3. --- 设 E 与 F 是局部凸空间，且 \(u \in \mathcal{L}(\mathrm{E};\mathrm{F})\)。设 \(\mathrm{E}_{1}\)（相应地，\(\mathrm{F}_{1}\)）是 E（相应地，F）的有限余维闭子空间。假设 u 把 \(\mathrm{E}_{1}\) 映到 \(\mathrm{F}_{1}\) 中，并用 \(u_{1} \in \mathcal{L}(\mathrm{E}_{1};\mathrm{F}_{1})\) 表示在 \(\mathrm{E}_{1}\) 上与 u 重合的映射。要使 u 是弗雷德霍姆算子，必须且只需 \(u_{1}\) 是。这时我们有

\[
\operatorname{ind} (u) - \operatorname{ind} \left(u _ {1}\right) = \operatorname{codim} _ {\mathrm{F}} \left(\mathrm{F} _ {1}\right) - \operatorname{codim} _ {\mathrm{E}} \left(\mathrm{E} _ {1}\right).\tag{5}
\]

设 \(i\colon \mathrm{E}_{1}\to \mathrm{E}\) 与 \(j\colon \mathrm{F}_{1}\to \mathrm{F}\) 是典范单射。它们是弗雷德霍姆映射，并且有

\[
\operatorname{ind} (i) = \operatorname{codim} _ {\mathrm{E}} (\mathrm{E} _ {1}), \quad \operatorname{ind} (j) = \operatorname{codim} _ {\mathrm{F}} (\mathrm{F} _ {1}).
\]

由于 \(j \circ u_{1} = u \circ i\)，可以看出 u 是弗雷德霍姆映射当且仅当 \(u_{1}\) 是（III, p. 41 的注记 3 和 4）。若是如此，我们有

\[
\operatorname{ind} (j) + \operatorname{ind} (u _ {1}) = \operatorname{ind} (j \circ u _ {1}) = \operatorname{ind} (u \circ i) = \operatorname{ind} (u) + \operatorname{ind} (i),
\]

由此得到公式 (5)。

命题 4. --- 设 E 与 F 是分离的局部凸空间，\(u: \mathrm{E} \to \mathrm{F}\) 是弗雷德霍姆算子，而 \(\widehat{u}: \widehat{\mathrm{E}} \to \widehat{\mathrm{F}}\) 是 u 到完备化上的延拓。那么 \(\widehat{u}\) 是弗雷德霍姆算子，并且有 \(\operatorname{Ker}(\widehat{u}) = \operatorname{Ker}(u)\) 与 \(\operatorname{ind}(\widehat{u}) = \operatorname{ind}(u)\)。

采用 III, p. 42 的 Prop. 2 的条件 (iv) 的记号。由于向量空间 \(E_{2}\) 与 \(F_{2}\) 是有限维的，它们是完备的。\(E_{1}\)（相应地，\(F_{1}\)）的完备化等同于 \(E_{1}\) 在 \(\widehat{E}\) 中的闭包（相应地，\(F_{1}\) 在 \(\widehat{F}\) 中的闭包），且 \(\widehat{E}\)（相应地，\(\widehat{F}\)）是 \(\widehat{E}_{1}\) 与 \(E_{2}\)（相应地，\(\widehat{F}_{1}\) 与 \(F_{2}\)）的拓扑直和。线性映射 \(\widehat{u}\) 通过限制定义了从 \(\widehat{E}_{1}\) 到 \(\widehat{F}_{1}\) 上的同构并在 \(E_{2}\) 上为零。命题于是由 loc. cit. 的蕴含关系 (iv)⇒(i) 得出。

